\section{Evaluation}
\label{sec:evaluation}

The evaluation is generated from repository scripts and processed CSVs; the paper-number audit rejects numeric claims that do not appear in generated evidence. The core full-RTL ledger reports \rcFullReplayPass/\rcFullReplayTotal\ compiler-backed record/replay PASS rows. The negative ledger rejects \rcNegativeReject\ replay-critical corruptions, reports \rcNegativeUnexpectedAccept\ unexpected accepts, and leaves \rcNegativeNA\ not-applicable rows for corruption classes outside the current commit-index harness.

\subsection{Replay, Scale, and Buffers}

Replay succeeds on the fixed compiler-backed suite in the full-RTL harness, but scaling is constrained by capsule capacity. The regenerated workload-scaling evidence reports 375/375 PASS rows with workload-aware v1 depths: 256 for smoke and short workloads, 1024 for medium, 4096 for long, and 16384 for stress. The buffer-sensitivity table shows replay success reaches 100.000000\% for each workload class at those selected depths, while the sticky overflow bit remains asserted on long-running fixed and no-failure rows that do not freeze at a property failure. The paper therefore treats buffer depth as a design parameter, not a hidden constant.

Capsule-size comparisons use the buffer depth required for the workload class under test. \texttt{capsule\_baseline\_summary.csv} reports median ReplayCapsule reductions versus the raw full-instruction-trace estimate of 59.52\% for smoke, 78.71\% for short, 83.18\% for medium, 64.86\% for long, and 22.84\% for stress. The same CSV includes a labeled \texttt{riscv\_etrace\_branch\_trace\_estimate} row; that estimate is smaller than ReplayCapsule-RV for these workloads, so the raw instruction-trace comparison is a weak baseline, not a state-of-the-art trace-compression claim.

The v2 evidence package adds compressed recorder configurations and two additional actuator/configuration firmware families. Those rows improve workload and benchmark coverage for the scoped model, but they do not turn the standalone replay-consume controller into an integrated full-core autonomous replay path.

\subsection{Cost and Mapping}

The mapped FPGA results are same-target bring-up/debug measurements. The board-level baseline and ReplayCapsule wrappers use the same \rcMappedTarget\ target and \rcMappedFlow\ flow. After gating unselected v2 recorder instances in v1 mapped builds, the representative mapped-scaling summary reports LUT overhead from 114.20\% to 221.23\%, FF overhead from 174.53\% to 646.43\%, BRAM overhead of 0.00\%, and Fmax change from -19.56\% to -9.98\%.

For the selected v2 minimal recorder profile, the same-target ECP5 measured row reports +8.26\% LUT, +3.77\% FF, +0.00\% BRAM, and -0.04\% Fmax. The core and hashed v2 rows remain diagnostic-rich comparison points, with +68.00\% and +67.69\% LUT overhead and +81.49\% FF overhead. These rows support the design direction that replay-critical boundary-event capture should be separated from optional on-chip diagnostics.

This overhead is acceptable only under the intended use case: temporary validation and post-failure replay instrumentation during hardware/software bring-up. It is not presented as a production always-on block. The current artifact does not include an ASIC or power flow; FPGA results should be read as feasibility and cost-shape evidence rather than signoff implementation data.

\ifrcfullartifactpaper
% Auto-generated by scripts/render_paper_tables.py from results/processed/*.csv.

\begin{table}[t]
\caption{Compiler-backed firmware benchmark suite. Source: \texttt{results/processed/firmware\_build.csv}.}
\label{tab:benchmarks}
\centering
\begin{tabular}{lll}
\toprule
Family & Variants & Failure mechanism \\
\midrule
Sensor threshold & failing, fixed, no\_failure\_edge & MMIO sensor threshold and deadline check \\
Interrupt race & failing, fixed & interrupt delivery inside a critical window \\
MMIO ordering & failing, fixed & unsafe command/output MMIO ordering \\
Stack corruption & failing, fixed & protected-store observation reaches checker \\
UART command & failing, fixed, no\_failure\_edge & unsafe UART command drives actuator write \\
Watchdog timeout & failing, fixed, no\_failure\_edge & missed heartbeat under long-running path \\
\bottomrule
\end{tabular}
\end{table}

\begin{table}[t]
\caption{Host-driven full RTL replay success. Source: \texttt{results/processed/full\_rtl\_replay.csv}.}
\label{tab:replay-success}
\centering
\begin{tabular}{lrrl}
\toprule
Evidence & Pass & Total & Firmware source \\
\midrule
Record/replay/signature match & 45 & 45 & compiler C \\
Compiler-backed rows & 45 & 45 & compiler C \\
\bottomrule
\end{tabular}
\end{table}

\begin{table}[t]
\caption{Full RTL corrupted-capsule behavior. Source: \texttt{results/processed/full\_rtl\_replay\_negative.csv}.}
\label{tab:negative}
\centering
\begin{tabular}{lr}
\toprule
Result & Rows \\
\midrule
REJECT & 10 \\
ACCEPT & 0 \\
NA & 2 \\
\bottomrule
\end{tabular}
\end{table}

\begin{table}[t]
\caption{Runtime and simulator-overhead summary. Source: \texttt{results/processed/runtime\_overhead\_summary.csv}. Simulator wall time is not a hardware timing claim.}
\label{tab:runtime}
\centering
\begin{tabular}{llrr}
\toprule
Config & Metric & Median & N \\
\midrule
baseline no recorder & sim wall time sec & 0.018194 & 45 \\
recorder present disabled & sim wall time sec & 0.024409 & 45 \\
recorder enabled & sim wall time sec & 0.024731 & 45 \\
recorder present disabled vs baseline no recorder & cycle overhead pct & 0.000000 & 45 \\
recorder enabled vs baseline no recorder & cycle overhead pct & 0.000000 & 45 \\
recorder enabled vs baseline no recorder & sim wall time overhead pct & 33.950000 & 45 \\
\bottomrule
\end{tabular}
\end{table}

\begin{table}[t]
\caption{Full-core mapped ECP5 overhead. Sources: \texttt{results/processed/mapped\_synthesis.csv} and \texttt{results/processed/mapped\_overhead.csv}.}
\label{tab:mapped-overhead}
\centering
\begin{tabular}{lrrrr}
\toprule
Design/metric & LUT & FF & BRAM & Fmax MHz \\
\midrule
Baseline board & 2814 & 883 & 6 & 63.47 \\
ReplayCapsule board & 6859 & 3901 & 6 & 50.70 \\
Overhead & 143.75\% & 341.79\% & 0.00\% & -20.12\% \\
\bottomrule
\end{tabular}
\end{table}

\begin{table}[t]
\caption{Scope and limitations.}
\label{tab:limitations}
\centering
\begin{tabular}{ll}
\toprule
Topic & Current scope \\
\midrule
Processor & single-hart RV32I/PicoRV32 \\
Nondeterminism & commit-indexed interrupt/MMIO boundary events \\
Replay engine & host-driven Verilator harness \\
Replay hardware & record-side recorder only; no replay-consume datapath \\
Memory system & no multicore, DMA, cache-coherence, or analog-device model \\
Optimization & fidelity and auditability prioritized over area minimization \\
\bottomrule
\end{tabular}
\end{table}

\fi
